\section{Conclusiones}\label{sec:conclusiones}
\begin{enumerate}
  \item Se formuló la asignación de vueltas a buses como un problema de máquinas paralelas idénticas con ventanas de inicio, almuerzo y tiempos de viaje $p_j(s_j)$ dependientes de la hora de salida, con objetivo lexicográfico: primero maximizar las vueltas atendidas y luego minimizar la carga máxima $\Lmax$.
  \item El tráfico se incorporó mediante un perfil de congestión por franjas que cumple la propiedad FIFO, evaluado en tres escenarios experimentales.
  \item Los resultados preliminares muestran que, sin tráfico, List Scheduling y LPT producen la misma solución, y que con tráfico LS atiende más vueltas y equilibra mejor la carga en las instancias evaluadas (en T1, 155 vueltas no atendidas frente a 293).
  \item En las 35 rutas evaluadas, el tráfico pesimista reduce la proporción de vueltas atendidas con LS de 98.2\,\% a 94.6\,\% y aumenta la conducción total en 13.4\,\%.
  \item La saturación de la flota ($\rho$) explica buena parte de las vueltas no atendidas; en RTU-22 la demanda de horas supera la capacidad de la flota.
  \item El prototipo permite verificar el modelo de extremo a extremo: genera las instancias con una regla general, ejecuta LS y LPT, valida todas las restricciones y produce resultados reproducibles, lo que muestra la viabilidad del proyecto.
\end{enumerate}
